% =====================================================================
\chapter{Introduction}
\label{ch:introduction}
% =====================================================================

\section{Motivation}
\label{sec:intro:motivation}

Embedded systems increasingly control safety- and time-critical
processes: motor controllers, medical devices, automotive electronic
control units, and industrial field buses all depend on firmware that
must react \emph{correctly} and \emph{on time}. Development and
verification of such firmware traditionally requires physical hardware,
which is expensive to scale, difficult to automate, and awkward to
instrument. \emph{Rehosting}---running unmodified firmware binaries
inside a full-system emulator---removes these obstacles and has become
a cornerstone of modern embedded workflows, from continuous-integration
testing~\cite{renode} to large-scale firmware security
analysis~\cite{fasano2021sok,wright2021challenges}.

Emulation frameworks such as Renode~\cite{renode} and
QEMU~\cite{bellard2005qemu} model the \emph{functional} contract of a
microcontroller: instruction semantics, memory maps, peripheral
registers, and interrupt controllers. Firmware that passes functional
tests in the emulator will, in most cases, execute the same logic on
real silicon. The \emph{temporal} contract, however, is a different
matter. Emulators advance virtual time in discrete scheduling quanta,
dispatch interrupts at quantum boundaries, execute instruction blocks at
an idealised rate, and omit analog-world phenomena such as crystal
oscillator drift, temperature-dependent jitter, and bus propagation
delay. As a result, a rehosted firmware image experiences a timing
environment that can differ from the physical one by orders of
magnitude---in either direction.

This discrepancy matters. A control loop that misses its deadline only
under real-world scheduling jitter will appear healthy in the emulator.
Conversely, quantum-induced noise in the emulator can trigger spurious
timeouts that never occur on hardware. In both cases the emulator's
verdict is misleading, undermining the very purpose of rehosted
testing. Despite this, the temporal accuracy of modern rehosting
frameworks has received little systematic, hardware-grounded
quantification.

\section{Problem Statement}
\label{sec:intro:problem}

This thesis addresses two questions:

\begin{enumerate}
  \item \textbf{RQ1 --- Quantification.} How large is the temporal
        deviation between an unmodified firmware image running on a
        physical Arm Cortex-M3 microcontroller and the same image
        rehosted in Renode, across representative RTOS timing
        primitives (sleeps, hardware timers, thread scheduling,
        interrupt service latency), and which mechanisms cause it?
  \item \textbf{RQ2 --- Mitigation.} Can the deviation be reduced by a
        lightweight, non-invasive extension layer---without modifying
        the emulator core, without patching the firmware, and without
        sacrificing the usability that makes rehosting attractive in
        the first place?
\end{enumerate}

\section{Contributions}
\label{sec:intro:contributions}

The thesis makes the following contributions:

\begin{itemize}
  \item \textbf{A measurement methodology and instrumented firmware}
        for comparing physical and virtual timing behaviour. A Zephyr
        RTOS firmware exercises ten timing scenarios and exposes each
        event as a GPIO edge, captured on hardware by a
        \SI{100}{\mega\sample\per\second} logic analyser and in Renode by a GPIO
        tracing hook (\cref{ch:problem}).
  \item \textbf{A taxonomy of five temporal deviation sources}
        (\srcS{1}--\srcS{5}), separating instruction-timing
        abstraction, interrupt delivery latency, clock drift and
        jitter, scheduler quantisation, and inter-node propagation
        delay (\cref{sec:problem:sources}).
  \item \textbf{The Temporal Fidelity Layer (TFL)}, three composable
        Renode peripheral models written in C\# that inject or monitor
        the missing temporal effects: \texttt{TFL\_ClockModel}
        (\srcS{3}), \texttt{TFL\_DeadlineMonitor} (\srcS{4}), and
        \texttt{TFL\_BusDelayModel} (\srcS{5}). TFL attaches purely via
        Renode platform descriptions and works with stock
        Renode~1.15.3 (\cref{ch:design}).
  \item \textbf{A statistical evaluation} of baseline Renode and
        TFL-enhanced Renode against hardware ground truth, including
        distribution comparisons, Mann--Whitney significance tests,
        a quantum-sensitivity study, and a workload-interference study
        (\cref{ch:evaluation}).
  \item \textbf{Open research artefacts}: firmware, peripheral models,
        platform files, capture data, and the full analysis pipeline
        are available in the accompanying repository
        (\cref{app:reproducibility}).
\end{itemize}

\section{Thesis Outline}
\label{sec:intro:outline}

\cref{ch:background} introduces the technical background on embedded
timing, the STM32F103 platform, Zephyr RTOS, and the Renode emulation
framework. \cref{ch:relatedwork} surveys related work on rehosting and
timing-accurate simulation. \cref{ch:problem} develops the measurement
methodology and the deviation-source taxonomy, and presents the
baseline characterisation of Renode against hardware.
\cref{ch:design} describes the design and implementation of the
Temporal Fidelity Layer. \cref{ch:evaluation} evaluates TFL
quantitatively. \cref{ch:conclusion} concludes and outlines future
work.
